% !TEX root = ../main.tex
%% Rozdział 5. Plan: 3–4 strony tekstu. Same fakty o badaniach; wyniki są w rozdziale 6.

\chapter{Metodyka badań}
\label{chap:metodyka}

Rozdział opisuje sposób przeprowadzenia badań generatora: postawione pytania, sprzęt i środowisko pomiarowe, przebieg pojedynczego pomiaru, plan eksperymentów, mierzone wielkości oraz sposób opracowania wyników. Same wyniki omawia \cref{chap:wyniki}.

\section{Pytania badawcze}
\label{sec:pytania}

Opis tematu pracy zakłada porównanie wariantów algorytmu \ac{wfc} pod względem czasu generacji, zużycia pamięci i kosztu rysowania gotowej mapy. Wynika z tego osiem pytań badawczych. Każde z nich ma przypisany eksperyment (\cref{tab:eksperymenty}).

\begin{enumerate}
  \item Jak sposób przechowywania możliwości komórek wpływa na czas generacji i zużycie pamięci? (E1)
  \item Która kolejność wyboru komórek daje najwięcej udanych prób i czy heurystyki \ac{mrv} oraz najmniejszej entropii sprawdzają się przy generacji zamków? (E2)
  \item Jak reakcja na sprzeczność wpływa na powodzenie i czas generacji przy rosnącym rozmiarze siatki? (E3)
  \item Jak czas generacji i pamięć rosną wraz z liczbą komórek siatki? (E4)
  \item Jak liczba kafelków w zestawie wpływa na czas i powodzenie generacji? (E5)
  \item Jak zmiany reguł i wag kafelków wpływają na zabudowę mapy i powodzenie generacji? (E6)
  \item Jak konsolidacja geometrii i atlas materiałów wpływają na koszt rysowania mapy? (E7)
  \item Czy miary restrykcyjności reguł pozwalają przewidzieć powodzenie generacji? (E8 i E9)
\end{enumerate}

Niezależnie od pytania każdy eksperyment sprawdza też poprawność map: każda gotowa mapa jest porównywana ze wszystkimi regułami zestawu kafelków (\cref{sec:metryki}).

\section{Środowisko pomiarowe}
\label{sec:srodowisko}

Wszystkie pomiary wykonano na jednym komputerze, opisanym w \cref{tab:sprzet}. Badania prowadzono w zbudowanej aplikacji \kol{Zamek.exe} dla 64-bitowego systemu Windows, a nie w edytorze Unity. Edytor wykonuje w tle własną pracę (odświeża okna, dodaje sprawdzenia bezpieczeństwa w kodzie silnika) i nie pozwala wyłączyć odśmiecania pamięci na czas pomiaru, co w zbudowanej aplikacji jest możliwe (\cref{sec:procedura}). Okno edytora i aplikacja korzystają z tej samej biblioteki \kol{WFC.Core} (\cref{sec:architektura}), więc wyniki dotyczą obu.

Aplikację zbudowano z backendem \ac{il2cpp} \cite{UnityIL2CPP} w konfiguracji Release: kod C\# jest zamieniany na C++, a ten kompilowany kompilatorem MSVC z włączonymi optymalizacjami. Wersje narzędzi podaje \cref{tab:technologie}. Wersja Development, która łączy się z oknem Unity Profiler \cite{UnityProfiler}, posłużyła wyłącznie do analizy działania programu i do zrzutów ekranu. Żadnej serii pomiarowej nie wykonano z podłączonym Profilerem, bo zbieranie przez niego danych wydłuża pracę programu. Liczby do tabel zbiera kod aplikacji.

Serie E1--E6, E8 i E9 mierzą wyłącznie generację, więc uruchomiono je bez okna, z argumentami \kol{-batchmode -nographics}. Seria E7 mierzy koszt rysowania, dlatego działała w oknie o rozmiarze 1600\,$\times$\,900 pikseli. Wszystkie argumenty uruchomienia zestawia \cref{tab:argumenty}. Serie E1--E6, E8 i E9 wykonano 10 i 11 listopada 2026 r., a serię E7 12 listopada 2026 r. Każdy wiersz plików wyników ma w kolumnie \kol{srodowisko} wartość \kol{IL2CPP}, czyli pochodzi z wersji Release.

\begin{table}
  \centering
  \footnotesize
  \caption{Środowisko pomiarowe}
  \label{tab:sprzet}
  \begin{tabular}{@{}ll@{}}
    \toprule
    Element & Wartość \\
    \midrule
    Procesor & AMD Ryzen 7 9800X3D, 8 rdzeni \\
    Pamięć operacyjna & 32 GB DDR5-6000, CL30 \\
    Karta graficzna & ASUS PRIME GeForce RTX 5080 OC Edition, 16 GB GDDR7 \\
    Sterownik karty graficznej & NVIDIA 610.88 \\
    System operacyjny & Windows 11 Pro \\
    Rozdzielczość okna & 1600\,$\times$\,900 pikseli, aplikacja w oknie \\
    Wersja Unity & 6000.6.2f1 \\
    Rodzaj aplikacji & IL2CPP, wersja Release \\
    \bottomrule
  \end{tabular}
\end{table}

\section{Procedura pomiaru}
\label{sec:procedura}

Eksperyment składa się z wariantów, czyli kombinacji ustawień generatora: rozwiązywacza, kolejności wyboru komórek, reakcji na sprzeczność, zestawu kafelków i rozmiaru siatki, a w E7 także sposobu rysowania i ustawienia łączenia rysowania. Warianty są wykonywane po kolei. Każdy zaczyna się od uruchomień rozgrzewkowych, których wyniki nie trafiają do pliku: domyślnie pięciu, w E7 trzech, a w E9 jednego. Rozgrzewka pochłania koszty pierwszego użycia kodu i danych, takie jak pierwsze alokacje pamięci, wczytanie modeli czy przygotowanie zadań kompilatora Burst \cite{UnityBurst}.

Po rozgrzewce następują uruchomienia mierzone, po jednym na każde ziarno generatora liczb losowych. Ziarna to kolejne liczby całkowite od 1, a ich liczbę w każdym eksperymencie podaje \cref{tab:eksperymenty}. Wszystkie warianty dostają ten sam zestaw ziaren, więc są porównywane na tych samych losowaniach. Generacja jest powtarzalna: to samo ziarno w tym samym wariancie zawsze daje tę samą mapę. Rozwiązywacze różnią się tylko sposobem przechowywania możliwości, więc z tego samego ziarna budują tę samą mapę. Potwierdza to E1: dla każdego ziarna wszystkie cztery rozwiązywacze wykonały tyle samo prób i decyzji, napotkały tyle samo sprzeczności, a gotowe mapy mają identyczną budowę.

Przed każdym uruchomieniem program dwukrotnie wywołuje odśmiecanie pamięci (\ac{gc}), czeka na zakończenie finalizatorów, a na czas generacji wyłącza odśmiecanie. Dzięki temu przerwa na odśmiecanie nie wpada w mierzony czas, a przyrost sterty po generacji jest równy temu, co zaalokował rozwiązywacz. Czas generacji mierzy sam rozwiązywacz zegarem \kol{Stopwatch}. Obejmuje on wszystkie próby danej mapy, razem ze startami od nowa i cofaniem decyzji. Osobno zapisywane są czasy wyboru komórki, propagacji i cofania. Sprawdzenie gotowej mapy z regułami i analiza jej budowy odbywają się już po zatrzymaniu zegara.

Generacja jednej mapy ma dwa limity: 1000 prób (w E9 30 prób) oraz, w E3 i E4, 600 s. Mapa, która przekroczy którykolwiek z nich, jest zapisywana jako nieudana (kolumna \kol{udane}), ale jej czas, pamięć i liczniki trafiają do pliku tak samo jak w pozostałych uruchomieniach. Mediany czasu wariantów z nieudanymi mapami obejmują więc także czas potrzebny do osiągnięcia limitu. W zebranych wynikach nieudane mapy wystąpiły tylko przy siatce \siatka{80}{80}{13} ze startem od nowa: rozwiązywacze \kol{Bitset} i \kol{Burst} zatrzymał limit 1000 prób (po 74--120~s), a rozwiązywacze \kol{Naive} i \kol{SupportCount} limit 600~s.

W E7 każda wygenerowana mapa jest stawiana w scenie. Pomiar rysowania zaczyna się dopiero pięć klatek później, bo liczniki silnika podają wartości z poprzedniej klatki, a pierwsze klatki po postawieniu mapy niosą jeszcze koszt jej stawiania. Następnie przez 120 klatek program odczytuje liczniki renderowania przez \kol{ProfilerRecorder} i zapisuje ich średnie na klatkę. Czas klatki to średni odstęp między kolejnymi klatkami (\kol{Time.unscaledDeltaTime}), a najdłuższa klatka to największy z tych odstępów. Ogólne liczniki wywołań rysowania nie działają w wersji Release, dlatego wywołania rysowania są sumą liczników osobnych dla każdej drogi rysowania w \ac{urp}, m.in. zwykłej, instancjonowanej, SRP Batchera i GPU Resident Drawera. Liczniki obejmują też cienie i panel aplikacji, co daje niewielki, stały narzut przy każdej mapie. Generacja w E7 używa rozwiązywacza \kol{Burst} z cofaniem decyzji, więc każda mapa powstaje w jednej próbie, a dla danego ziarna wszystkie dwanaście ustawień rysowania pokazuje tę samą mapę.

Każde mierzone uruchomienie zapisuje się jako jeden wiersz pliku \ac{csv}. Po zakończeniu serii aplikacja zapisuje obok podsumowanie wariantów (plik \kol{_warianty.csv}), a w eksperymentach z różnymi regułami także korelacje (plik \kol{_korelacje.csv}). Kolumny plików opisują \cref{tab:kolumny-csv,tab:pliki-podsumowania}.

\section{Plan eksperymentów}
\label{sec:plan}

Ustawienia wszystkich eksperymentów zestawia \cref{tab:eksperymenty}, a rozmiary siatki \cref{tab:rozmiary}. Eksperymenty, które nie badają danego ustawienia, używają zestawu \kol{Pelny}, siatki \siatka{20}{20}{13}, kolejności \kol{LevelFirst} i startu od nowa (\kol{Restart}).
\begin{itemize}
  \item E1 porównuje cztery rozwiązywacze (\cref{tab:rozwiazywacze}).
  \item E2 porównuje pięć kolejności wyboru komórek (\cref{tab:heurystyki}).
  \item E3 porównuje trzy reakcje na sprzeczność (\cref{tab:reakcje}) na czterech rozmiarach siatki.
  \item E4 mierzy czas i pamięć czterech rozwiązywaczy na tych samych czterech rozmiarach. W E3 i E4 limit 600~s na mapę nie pozwala najwolniejszym wariantom zablokować serii.
  \item E5 porównuje cztery zestawy kafelków o rosnącej liczbie stanów (\cref{tab:zestawy}).
  \item E6 porównuje zestaw \kol{Pelny} z czterema jego wariantami reguł (\cref{tab:warianty-regul}).
  \item E7 porównuje cztery sposoby rysowania mapy (\cref{tab:rysowanie}) przy trzech ustawieniach łączenia rysowania w \ac{urp} (\cref{tab:laczenie}) i czterech rozmiarach siatki.
  \item E8 zestawia restrykcyjność reguł ośmiu zestawów i wariantów z E5 i E6 z odsetkiem udanych prób.
  \item E9 zaostrza reguły zestawu \kol{Pelny}: program usuwa losowo 0,2, 0,5, 1, 2, 3, 5, 8 albo 12\% dozwolonych par sąsiadów, po dziesięć losowań na poziom. Razem z regułami niezmienionymi daje to 81 wariantów. Limit 30 prób na mapę skraca serię przy regułach, które rzadko dają mapę.
\end{itemize}

Łącznie wykonano 8460 mierzonych uruchomień (\cref{tab:naruszenia}).

\begin{table}
  \centering
  \footnotesize
  \caption{Plan eksperymentów}
  \label{tab:eksperymenty}
  \begin{tabularx}{\linewidth}{@{}lLrrL@{}}
    \toprule
    Nr & Co się zmienia & Warianty & Ziarna & Ustawienia stałe \\
    \midrule
    E1 & rozwiązywacz (4) & 4 & 100 & \kol{Pelny}, \siatka{20}{20}{13}, \kol{LevelFirst}, \kol{Restart} \\
    E2 & heurystyka (5) & 5 & 100 & \kol{Pelny}, \siatka{20}{20}{13}, \kol{Bitset}, \kol{Restart} \\
    E3 & reakcja (3) $\times$ rozmiar (4) & 12 & 100 & \kol{Pelny}, \kol{Bitset}, \kol{LevelFirst}, limit 600 s na mapę \\
    E4 & rozwiązywacz (4) $\times$ rozmiar (4) & 16 & 100 & \kol{Pelny}, \kol{LevelFirst}, \kol{Restart}, limit 600 s na mapę \\
    E5 & zestaw (4) & 4 & 100 & \siatka{20}{20}{13}, \kol{Bitset}, \kol{LevelFirst}, \kol{Restart} \\
    E6 & \kol{Pelny} i warianty reguł (4) & 5 & 100 & \siatka{20}{20}{13}, \kol{Bitset}, \kol{LevelFirst}, \kol{Restart} \\
    E7 & rysowanie (4) $\times$ łączenie (3) $\times$ rozmiar (4) & 48 & 30 & \kol{Pelny}, \kol{Burst}, \kol{LevelFirst}, \kol{Backtrack}, rozgrzewka 3, 120 klatek na mapę, tylko w oknie \\
    E8 & 4 zestawy i 4 warianty reguł & 8 & 100 & \siatka{20}{20}{13}, \kol{Burst}, \kol{LevelFirst}, \kol{Restart} \\
    E9 & zaostrzenie reguł (9 poziomów) $\times$ 10 losowań & 81 & 20 & \kol{Pelny}, \siatka{20}{20}{13}, \kol{Burst}, \kol{LevelFirst}, \kol{Restart}, rozgrzewka 1, limit 30 prób \\
    \bottomrule
  \end{tabularx}
\end{table}

\begin{table}
  \centering
  \footnotesize
  \caption{Rozmiary siatki}
  \label{tab:rozmiary}
  \begin{tabular}{@{}lrl@{}}
    \toprule
    E\,$\times$\,N\,$\times$\,piętra & Komórki & Wymiar mapy [m] \\
    \midrule
    \siatka{10}{10}{8} & 800 & \siatka{30}{30}{24} \\
    \siatka{20}{20}{13} & 5200 & \siatka{60}{60}{39} \\
    \siatka{40}{40}{13} & 20\,800 & \siatka{120}{120}{39} \\
    \siatka{80}{80}{13} & 83\,200 & \siatka{240}{240}{39} \\
    \bottomrule
  \end{tabular}
\end{table}

\section{Metryki}
\label{sec:metryki}

Główne mierzone wielkości zestawia \cref{tab:metryki}, a pełną listę kolumn pliku wyników \cref{tab:kolumny-csv}. Poniżej opisano te, których znaczenie nie wynika wprost z nazwy.

\begin{table}
  \centering
  \footnotesize
  \caption{Główne metryki}
  \label{tab:metryki}
  \begin{tabular}{@{}lll@{}}
    \toprule
    Metryka & Kolumna pliku wyników & Jednostka \\
    \midrule
    Czas generacji & \kol{czas_ms} & ms \\
    Czas wyboru komórki & \kol{czas_wybor_ms} & ms \\
    Czas propagacji & \kol{czas_propagacja_ms} & ms \\
    Czas cofania & \kol{czas_cofanie_ms} & ms \\
    Próby i sprzeczności & \kol{proby}, \kol{sprzecznosci} & liczba \\
    Naruszenia zasad & \kol{naruszenia} & liczba \\
    Pamięć szacowana & \kol{pamiec_szczytowa_mb} & MB \\
    Pamięć rzeczywista & \kol{pamiec_rzeczywista_mb} & MB \\
    Wywołania rysowania & \kol{wywolania_rysowania} & liczba \\
    Narysowane obiekty & \kol{narysowane_obiekty} & liczba \\
    Zmiany stanu SetPass & \kol{setpass} & liczba \\
    Czas klatki & \kol{czas_klatki_ms} & ms \\
    Najdłuższa klatka & \kol{najdluzsza_klatka_ms} & ms \\
    Czas stawiania mapy & \kol{czas_stawiania_ms} & ms \\
    Pamięć silnika & \kol{pamiec_unity_mb} & MB \\
    \bottomrule
  \end{tabular}
\end{table}

\paragraph{Pamięć.} Pamięć szacowana to największy rozmiar struktur danych rozwiązywacza, liczony z rozmiarów jego tablic. Pamięć rzeczywista to przyrost sterty po generacji i pełnym odśmieceniu, zmierzony, gdy rozwiązywacz jeszcze istnieje, czyli pamięć, którą rozwiązywacz naprawdę zajmuje. Tabele wyników podają pamięć rzeczywistą. Pamięć silnika (E7) to cała pamięć zaalokowana przez Unity według licznika silnika.

\paragraph{Udane próby.} Odsetek udanych prób wariantu wynosi
\begin{equation}
  u = 100 \cdot \frac{M}{A} ,
  \label{eq:udane-proby}
\end{equation}
gdzie $M$ to liczba udanych map wariantu, a $A$ to liczba wszystkich prób wszystkich jego map. Przy starcie od nowa każdy start jest osobną próbą, więc $u$ mówi, jak często jedno przejście algorytmu dochodzi do końca bez sprzeczności.

\paragraph{Zabudowa i mapy zdegenerowane.} Zabudowa to odsetek komórek parteru zajętych przez budowle, uśredniony po udanych mapach wariantu. Wariant, w którym ta średnia jest mniejsza niż 1\% albo większa niż 90\%, ma mapy zdegenerowane: reguły usunęły z mapy budynki albo teren. Mapa powstaje wtedy łatwo, ale nie jest zamkiem, dlatego korelacje w E8 i E9 są liczone także bez takich wariantów.

\paragraph{Różnorodność map.} Każda udana mapa jest porównywana z ostatnią udaną mapą poprzedniego ziarna w tym samym wariancie. Miarą jest odsetek komórek z innym kafelkiem, osobno dla całej mapy, parteru i zabudowy.

\paragraph{Naruszenia zasad.} Każda udana mapa jest sprawdzana ze wszystkimi regułami zestawu kafelków niezależnie od rozwiązywacza. Kolumna \kol{naruszenia} podaje liczbę miejsc, w których mapa łamie reguły.

\paragraph{Propagacja.} Możliwości komórek znikają na dwa sposoby: przy decyzji, gdy wybrana komórka zostaje ustalona na jeden stan, i w propagacji, która usuwa stany niezgodne z sąsiadami. Udział propagacji $q$ i liczbę możliwości usuniętych przez propagację na jedną decyzję $e$ opisuje
\begin{equation}
  q = \frac{r_p}{r_p + r_d} , \qquad e = \frac{r_p}{D} ,
  \label{eq:propagacja}
\end{equation}
gdzie $r_p$ i $r_d$ to możliwości usunięte odpowiednio przez propagację (\kol{usuniete_propagacja}) i przez decyzje (\kol{usuniete_decyzje}), a $D$ to liczba decyzji (\kol{decyzje}). W tabelach $q$ jest podawane w procentach.

\paragraph{Obiekty na wywołanie rysowania.} Skuteczność łączenia rysowania w E7 opisuje liczba obiektów narysowanych jednym wywołaniem
\begin{equation}
  b = \frac{O}{V} ,
  \label{eq:obiekty-na-wywolanie}
\end{equation}
gdzie $O$ to narysowane obiekty (\kol{narysowane_obiekty}), a $V$ to wywołania rysowania (\kol{wywolania_rysowania}).

\paragraph{Restrykcyjność reguł.} W E8 i E9 porównywane są trzy miary restrykcyjności reguł (kod: \kol{RuleRestrictiveness.cs}). Niech $m$ oznacza liczbę stanów, $P$ liczbę dozwolonych trójek (stan, kierunek, sąsiad), $N_k(s)$ zbiór dozwolonych sąsiadów stanu $s$ w kierunku $k$, $w_s$ wagę stanu $s$, a $W$ sumę wag. Gęstość par to odsetek dozwolonych par sąsiadów wśród wszystkich możliwych w sześciu kierunkach:
\begin{equation}
  d = \frac{P}{6 m^2} .
  \label{eq:gestosc}
\end{equation}
Gęstość ważona uwzględnia, jak często stany pojawiają się zgodnie z wagami:
\begin{equation}
  d_w = \frac{1}{6} \sum_{k=1}^{6} \sum_{s=1}^{m} \frac{w_s}{W} \sum_{t \in N_k(s)} \frac{w_t}{W} .
  \label{eq:gestosc-wazona}
\end{equation}
Trzecią miarą jest współczynnik $\kappa$ \cite{Gent1996}, opisujący, jak bardzo ograniczony jest problem spełniania ograniczeń:
\begin{equation}
  \kappa = -\frac{c}{n} \log_m d ,
  \label{eq:kappa}
\end{equation}
gdzie $c/n$ to liczba ograniczeń na jedną zmienną. W siatce trójwymiarowej każda komórka ma sześciu sąsiadów, a każde ograniczenie łączy dwie komórki, więc przyjęto $c/n = 3$. Im większe $\kappa$, tym ostrzejsze reguły, dlatego korelacja $\kappa$ z powodzeniem ma znak przeciwny niż korelacja gęstości. Tabele podają $d$ i $d_w$ w procentach.

\section{Analiza statystyczna}
\label{sec:statystyka}

Czasy generacji mają rozkład silnie prawoskośny: większość map powstaje w pierwszej próbie, a nieliczne potrzebują wielu startów od nowa (\cref{fig:e1-czas}). Dlatego wyniki są podawane jako mediana oraz kwartyle $Q_1$ i $Q_3$, a nie jako średnia i odchylenie standardowe. Rozrzut opisuje rozstęp międzykwartylowy
\begin{equation}
  \mathrm{IQR} = Q_3 - Q_1 .
  \label{eq:iqr}
\end{equation}
Kwantyle są liczone z interpolacją liniową między sąsiednimi wartościami, tak samo w aplikacji (\kol{Statistics.cs}) i w skryptach analizy. Wyjątkiem są liczby elementów zamku (wież, donżonów, mostów, schodów), ich wysokości i różnorodność map. Te wielkości podawane są jako średnie na mapę, bo mediany liczb całkowitych słabo rozróżniają warianty. Podpisy tabel zaznaczają, które wartości są średnimi, a które sumami.

Różnice czasu generacji między dwoma wariantami sprawdzono dwustronnym testem Manna-Whitneya \cite{MannWhitney1947}, który nie zakłada rozkładu normalnego. Przyjęto poziom istotności 0,05. Test zastosowano do par rozwiązywaczy w E1 (\cref{tab:e1-mann-whitney}).

Związek miar restrykcyjności z odsetkiem udanych prób opisuje współczynnik korelacji rang Spearmana \cite{Spearman1904}
\begin{equation}
  \rho = \mathrm{corr}\bigl(\mathrm{rg}(X), \mathrm{rg}(Y)\bigr) ,
  \label{eq:spearman}
\end{equation}
czyli współczynnik korelacji liniowej rang $\mathrm{rg}$ obu wielkości. Wartości równe dostają średnią rangę. Korelację policzono na wariantach (8 w E8 i 81 w E9), także bez wariantów z mapami zdegenerowanymi.

Tempo wzrostu czasu i pamięci wraz z liczbą komórek $n$ opisuje nachylenie $a$ prostej
\begin{equation}
  \log t = a \log n + b ,
  \label{eq:nachylenie}
\end{equation}
dopasowanej metodą najmniejszych kwadratów do median czasu albo pamięci $t$. Wartość $a \approx 1$ oznacza wzrost liniowy, a $a \approx 2$ kwadratowy. Nachylenie czasu dopasowano tylko do trzech mniejszych siatek, bo przy \siatka{80}{80}{13} część map nie powstaje, a mediana zależy wtedy od limitów, a nie od samego rozwiązywacza.

Pliki wyników opracowano w języku Python z bibliotekami pandas, SciPy i matplotlib (skrypty \kol{Analiza/analiza_wynikow.py} i \kol{Analiza/wykresy.py}). Skrypt wykresów liczy wartości na wariant tak samo jak aplikacja i porównuje je z plikiem \kol{_warianty.csv}. W zebranych wynikach nie wykazał żadnych różnic.
